% !TEX root = task-3-17-6.tex
% Editable task bank. Append new BankTask blocks at the end of the document.
\documentclass[12pt,a4paper]{article}
\usepackage{physics-textbook}
\usepackage{graphicx,needspace}
\usetikzlibrary{decorations.pathmorphing}
\newif\ifShowSolutions
\ifdefined\TutorEdition
  \ShowSolutionstrue
\else
  \ShowSolutionsfalse % Change to \ShowSolutionstrue to display solutions.
\fi
\newcounter{banktask}
\newsavebox{\BankTaskBox}
\newcommand{\BankTaskContent}[3]{%
  \noindent\textbf{\thebanktask.}\quad #2\par
  \ifShowSolutions
    \Needspace{50mm}% Keep the solution heading with its opening text or diagram.
    #3\par
    \ifdefined\TutorEdition
      \ifcsname TutorSource#1\endcsname
        \nopagebreak\smallskip{\small\csname TutorSource#1\endcsname\par}%
      \fi
    \fi
  \fi
}
\newcommand{\BankTask}[3]{%
  \refstepcounter{banktask}\par\medskip
  \sbox{\BankTaskBox}{\begin{minipage}{\linewidth}%
    \BankTaskContent{#1}{#2}{#3}\end{minipage}}%
  \ifdim\dimexpr\ht\BankTaskBox+\dp\BankTaskBox\relax>\dimexpr\textheight-20mm\relax
    \Needspace{35mm}\BankTaskContent{#1}{#2}{#3}%
  \else
    \Needspace{\dimexpr\ht\BankTaskBox+\dp\BankTaskBox\relax}%
    \noindent\usebox{\BankTaskBox}\par
  \fi
  \medskip
}
\begin{document}
\BookSubsection{3.17.6}{Магнитное поле катушки с током}


% Source 063141: https://ege.fipi.ru/bank/questions.php?search=1&proj=BA1F39653304A5B041B656915DC36B38&qid=063141
\BankTask{063141}{
Катушка №~1 включена в электрическую цепь, состоящую из источника постоянного напряжения и реостата. Катушка №~2 помещена внутрь катушки №~1, и её обмотка замкнута. Вид с торца катушек представлен на рисунке.

Ползунок реостата перемещают влево.
\begin{enumerate}
\item Сила тока в катушке №~1 увеличивается.
\item Модуль вектора индукции магнитного поля, созданного катушкой №~1, увеличивается.
\item Модуль магнитного потока, пронизывающего катушку №~2, уменьшается.
\item Вектор магнитной индукции магнитного поля, созданного катушкой №~2 в её центре, направлен от наблюдателя.
\item В катушке №~2 индукционный ток направлен по часовой стрелке.
\end{enumerate}
\begin{center}
\begin{tikzpicture}[book diagram,x=0.85cm,y=0.85cm]
\draw (1.65,2.25) ++(-104:0.9) arc[start angle=-104,end angle=-436,radius=0.9];
\draw (1.65,2.25) circle[radius=0.5];
\node at (1.65,2.37) {№ 2};
\node at (0.36,1.98) {№ 1};
\draw (1.432,1.377)--(0.2,1.377)--(0.2,0.38)--(0.58,0.38);
\draw (1.868,1.377)--(3.14,1.377)--(3.14,0.38)--(2.35,0.38);
\draw (0.58,0.18) rectangle (1.36,0.58);
\draw (1.36,0.38)--(2.18,0.38);
\draw (1.77,0.38)--(1.77,1.00)--(0.98,1.00);
\draw[-{book tip}] (0.98,1.00)--(0.98,0.60);
\draw (2.18,-0.04)--(2.18,0.80);
\draw (2.35,0.16)--(2.35,0.60);
\node at (1.88,-0.05) {$+$};
\node at (2.62,-0.05) {$-$};
\end{tikzpicture}
\end{center}
}{
\smallskip\noindent\textbf{Решение от GPT}\par
При перемещении ползунка влево сопротивление включённой части реостата уменьшается, поэтому ток в катушке №~1 и создаваемая им индукция возрастают: утверждения 1 и 2 верны. Модуль магнитного потока через катушку №~2 увеличивается, поэтому утверждение 3 неверно. Ток от положительного полюса источника идёт по катушке №~1 по часовой стрелке, создавая поле от наблюдателя. По правилу Ленца поле индукционного тока катушки №~2 направлено к наблюдателю, а сам ток идёт против часовой стрелки: утверждения 4 и 5 неверны. Ответ: 1, 2.
}

\end{document}
